\section{Introduction}
\label{sec:intro}

Security operations centers receive far more alerts than analysts can investigate, and most of those alerts turn out to be benign~\citep{nodoze}. LLM agents promise relief: they query logs, read the results, and decide what caused an alert~\citep{excytin,clouseau}. Organizations that may not send telemetry to a hosted model can run such an agent on a small open-weight model inside their own network. That keeps the data in-house, but it puts the whole investigation in the hands of a weaker model.

An investigation consists of two different jobs. The first is \emph{reading}: turning what a probe returns, often raw log text, into evidence such as ``the failures come from many residential addresses''. The second is \emph{deciding}: which probe to run next, whether the evidence already suffices, and which verdict to accept. Current agents give both jobs to the model~\citep{react}. We find that small models are poor at the second and indispensable for the first. Qwen2.5-7B choosing its own probes resolves 8 to 13\% of cases while spending nearly its whole budget, and Llama-3.1-8B requests another probe in all 3{,}389 of its decisions, even when its own ledger shows a posterior above 0.999. Yet when a logging change rewords the log lines, a fixed parser reads none of them, while the same 7B model reads them correctly.

Our key idea is to divide the jobs accordingly: let the model read, but not decide. We present \hesp{}, a controller for local triage agents. Given the alert's candidate causes and a catalogue of read-only probes, \hesp{} keeps a Bayesian ledger over the causes, runs the probe with the highest expected information gain per cost, ends the investigation itself once one cause is supported by current evidence, and accepts a verdict only if the ledger backs it. The model proposes next steps and, when probes return raw text, parses it into the outcomes the ledger uses.

The division raises three questions, and \hesp{} answers each with a rule. First, does taking the decisions from the model lose cases? It does not: the controller alone resolves every case in two of our three settings, and no model inside it, from 7B to 72B, resolves more; the strong models' only advantage, waiting for a confirming observation, is reproduced by a stopping rule. Second, is the model needed at all? Only for reading: with structured observations the controller needs no model, but under a changed log format every case escalates without one. Third, is a model reader safe? It is not: an attacker who writes a consistent benign story into every log the agent reads makes the model adopt it everywhere, which defeats even corroboration across independent sources. \hesp{} therefore records which parser produced each observation and lets model-parsed observations confirm attacks but never benign verdicts. Two further rules protect the decisions themselves: a finish guard rejects verdicts that the ledger does not support, which stops instructions injected into logs, and benign verdicts must be corroborated by two independent source groups, which stops a single forged source.

We evaluate \hesp{} on three simulated triage settings, one derived from public Sigma detection rules, with five open-weight models from 7B to 72B and three classes of attack, including adaptive attacks written against \hesp{} itself. The settings make two assumptions that define the scope of our results: the list of candidate causes is given in advance, and each cause leaves at least one observation that singles it out. Under these assumptions the procedure, not the model, does the deciding. Whether a model contributes to deciding when the candidates must be proposed during the investigation, or when evidence only adds up in combination, is a question our settings cannot answer.

In summary, we make the following contributions:
\begin{itemize}
  \item A measurement of where small local models fail in alert triage: they fail at choosing probes and at stopping, for different reasons in models of similar size, while a controller without any model resolves as many cases as every model we tested (\S\ref{sec:decide}).
  \item Evidence that the model's necessary role is reading: under a changed log format, a fixed parser escalates every case and a 7B reader resolves most of them (\S\ref{sec:read}).
  \item \hesp{}, a controller that lets a local LLM read but not decide, with three rules that keep attackers away from the verdict: reader trust, a finish guard, and source-diverse corroboration, evaluated against injected instructions, forged sources, and adaptive attacks on the reader (\S\ref{sec:approach}, \S\ref{sec:rawlog}, \S\ref{sec:injection}).\footnote{Code, prediction tables, outcome files, and episode journals are available at \repourl.}
\end{itemize}
